% !PRINTED 698
\weberpage{698}
Nach (7) und \S\,192 ist
\weq{9}{(O, B) \geqq (O, B') - m,
}
und nach (6), angewandt auf $A'$:
\[
(O, A) = (O, A') = a - m - p + 1 + (O, B'),
\]
also nach (8):
\[
(O, A) - \tfrac{1}{2} a \leqq (O, B) - \tfrac{1}{2} b.
\]

Da wir aber $A$ und $B$ vertauschen können, so ergibt sich
hieraus
\[
(O, A) - \tfrac{1}{2} a = (O, B) - \tfrac{1}{2} b,
\]
und wenn man
\[
(O, B) = (A, W),
\]
\[
\tfrac{1}{2}(a - b) = a - p + 1
\]
setzt:
\[
(O, A) = a - p + 1 + (A, W)
\]
in genauer Übereinstimmung mit (6).

Um auch für uneigentliche Klassen zweiter Gattung den
entsprechenden Satz abzuleiten, sei
\weq{10}{A = \mathfrak{M} A'}
eine uneigentliche Klasse zweiter Gattung vom Teiler $\mathfrak{M}$.

In einer eigentlichen Klasse zweiter Gattung, $C$, deren
Dimension größer ist als die Ordnung $a$ von $A$, gibt es nach
\S\,192, 5. immer Polygone, die durch ein Polygon der Klasse $A$
teilbar sind, und $C$ ist also auch durch $A$ teilbar.

Aus \S\,192, (15) folgt:
\[
(O, A) \geqq (O, C) - c + a,
\]
und aus (3):
\[
(O, C) = c - p + 1,
\]
also
\weq{11}{(O, A) \geqq a - p + 1.
}

Nun ist $A'$ eine eigentliche Klasse von derselben Dimension
wie $A$, also, wenn $m$ die Ordnung von $\mathfrak{M}$ ist, nach (6):
\weq{12}{(O, A) = (O, A') = a - m - p + 1 + (A', W),}
also nach (11):
\weq{13}{(A', W) \geqq m.
}

\gesperrt{Daraus folgt, da\ss\ $A'$ von der ersten Gattung sein
mu\ss.}

Ist $W = A'B'$, so folgt aus (13):
\weq{14}{(A', W) = (O, B') \geqq m.
}
